\section{总结与展望}\label{ch:conclusion}

\subsection{工作总结}\label{sec:summary}

本文围绕基于视觉语言模型的多曝光图像融合方法及其工程应用展开研究，主要完成了以下工作：

第一，复现并适配了一种基于视觉语言模型的图像融合方法FILM（Image Fusion via Vision-Language Model）。该方法利用预提取的文本语义特征，通过跨注意力机制实现文本语义对视觉特征融合过程的引导。本文结合代码实现分析了包含文本特征预处理、视觉-文本跨注意力融合块、三级级联融合结构在内的完整网络架构，并将其适配到多曝光图像融合任务中。在MEFB数据集上的实验结果表明，FILM方法在SD、SF、AG、VIF和Qabf五项评价指标上均优于传统融合方法，EN指标与均值融合持平但低于部分传统方法，验证了文本语义引导机制在多曝光融合任务中的有效性。

第二，设计并实现了一个基于Web的交互式多曝光图像融合系统。系统采用前后端分离的架构，前端基于Vue 3框架构建单页应用，提供图像上传、算法选择、融合结果可视化和质量评估等交互功能；后端基于FastAPI框架构建RESTful API服务，负责加载预训练的FILM模型并执行融合推理计算。系统支持AI融合（FILM模型）和传统融合算法（均值融合、像素最大值融合、Mertens融合）两种模式，满足了不同用户的需求。功能测试、性能测试和兼容性测试结果表明，系统各项功能运行正常，性能满足预期要求。

\subsection{不足与展望}\label{sec:future}

尽管本文在基于视觉语言模型的图像融合方法适配和系统实现方面取得了一定的成果，但仍存在一些不足之处，需要在未来的研究中进一步改进和完善：

（1）\textbf{文本语义来源的局限性}：当前系统端采用从MEFB测试数据中读取的固定文本特征作为语义提示，尚未针对用户上传的每一组图像实时生成文本描述。未来可以接入在线视觉语言模型，根据输入图像自动生成或更新文本特征，也可以探索人工标注的细粒度语义标签、多层次语义描述和用户自定义文本提示，进一步提升语义引导的针对性。

（2）\textbf{融合任务的泛化能力}：本文主要在多曝光图像融合任务上进行了实验验证和系统开发，对于红外-可见光融合、医学图像融合、多聚焦图像融合等其他融合任务的适用性尚未充分验证。未来可以将FILM方法扩展到更多的融合任务中，验证其泛化能力。

（3）\textbf{系统性能的优化}：当前系统的AI融合处理时间约为2.5秒，对于实时性要求较高的应用场景仍存在一定延迟。未来可以通过模型压缩、量化、知识蒸馏等技术手段对FILM模型进行轻量化处理，提高推理速度。同时，可以考虑部署到云端服务器，为用户提供更高效的在线融合服务。

（4）\textbf{用户交互体验的提升}：当前系统主要提供了基础的融合功能，用户可以调节的参数较少。未来可以增加更多的交互选项，如融合强度调节、区域选择融合、自定义文本描述输入等，使用户能够更精细地控制融合过程。同时，可以引入批量处理功能，提高用户的工作效率。

（5）\textbf{评价指标的完善}：当前系统仅实现了EN、SD、SF、AG四项客观评价指标，评价维度相对有限。未来可以引入更多的评价指标，如结构相似性（SSIM）、感知图像质量（LPIPS）、弗雷歇距离（FID）等，为用户提供更全面的融合质量评估。同时，可以引入主观评价机制，收集用户的主观满意度评分，建立更完善的评价体系。
